\documentclass[12pt,a4paper]{ctexart}

\usepackage[a4paper,margin=1.8cm]{geometry}
\usepackage{amsmath,amssymb,bm}
\usepackage{booktabs,longtable,tabularx,array,multirow}
\usepackage{tikz}
\usetikzlibrary{arrows.meta,calc,positioning,angles,quotes,mindmap,shadows,decorations.pathreplacing}
\usepackage[most]{tcolorbox}
\usepackage{enumitem}
\usepackage{multicol}
\usepackage{hyperref}
\usepackage{xcolor}
\usepackage{fancyhdr}
\setlength{\headheight}{15pt} % ctex 12pt + 中文页眉行距需要，避免 fancyhdr 报 headheight 过小
\addtolength{\topmargin}{-3pt} % 与上句配套，保持版芯垂直位置

\hypersetup{
  colorlinks=true,
  linkcolor=black,
  urlcolor=blue
}

\pagestyle{fancy}
\fancyhf{}
\fancyhead[L]{初中物理·简单机械拔高复习讲义}
\fancyhead[R]{第一版初稿}
\fancyfoot[C]{\thepage}

\setlength{\parindent}{2em}
\setlength{\parskip}{0.25em}
\setlist[enumerate]{itemsep=0.35em, topsep=0.35em}

\newtcolorbox{modelbox}[1]{
  colback=blue!3,
  colframe=blue!55!black,
  title=#1,
  sharp corners,
  boxrule=0.8pt
}

\newtcolorbox{pitfallbox}[1]{
  colback=red!3,
  colframe=red!55!black,
  title=#1,
  sharp corners,
  boxrule=0.8pt
}

\newcommand{\Ans}[1]{\par\noindent\textbf{答案：}#1\par}
\newcommand{\Sol}[1]{\par\noindent\textbf{解析：}#1\par}
\newcommand{\Tip}[1]{\par\noindent\textbf{点拨：}#1\par}

\begin{document}

\begin{center}
{\LARGE \textbf{初中物理·简单机械}}\\[0.4em]
{\Large \textbf{中考—竞赛衔接拔高复习讲义（第一版）}}\\[0.8em]
适用对象：初三学生\\
难度配比：基础巩固 20\%，中考提升 50\%，竞赛拔高 30\%\\
\end{center}

\noindent\textbf{使用约定：}
本讲义默认取 $g=10\,\mathrm{N/kg}$，水的密度取 $\rho_{\text{水}}=1.0\times 10^3\,\mathrm{kg/m^3}$。
课堂使用时建议先讲“模型”，再做“变式”，最后做“综合”。

\tableofcontents

\section*{总览与学习策略}
\addcontentsline{toc}{section}{总览与学习策略}

\begin{center}
\small
\begin{tabularx}{\textwidth}{>{\raggedright\arraybackslash}p{3.1cm} X c c c}
\toprule
章节 & 重点内容 & 基础题 & 提升题 & 拔高题\\
\midrule
第一讲 杠杆模型 & 五要素、力臂、平衡、最小力、动态、与浮力综合 & 3 & 4 & 3 \\
第二讲 滑轮与滑轮组 & 定/动滑轮、绳股数、力—距离—速度、摩擦参与 & 3 & 4 & 3 \\
第三讲 功、功率与机械效率 & 有用功、总功、效率模型、实验探究、功率综合 & 3 & 4 & 3 \\
第四讲 中考压轴与竞赛综合 & 图像题、极值、多过程、杠杆+滑轮+摩擦/浮力综合 & 3 & 4 & 3 \\
\bottomrule
\end{tabularx}
\end{center}

\begin{modelbox}{全书总策略}
\begin{enumerate}[label=\arabic*.]
\item 先定物理模型：杠杆、滑轮、机械效率、综合多过程。
\item 再找主方程：杠杆写力矩平衡；滑轮先找绳股数 $n$；效率先分有用功与总功。
\item 再找“隐藏量”：表面是力，实质往往考力臂；表面是效率，实质往往考功与位移；表面是图像，实质往往考线性关系。
\item 一律做单位检查：长度统一到米，力统一到牛，功统一到焦。
\end{enumerate}
\end{modelbox}

\begin{center}
\begin{tikzpicture}[mindmap,grow cyclic,
every node/.style={concept, circular drop shadow, text=black, minimum size=1.6cm},
concept color=blue!25,
level 1/.append style={level distance=4.0cm,sibling angle=90},
level 2/.append style={level distance=2.6cm,sibling angle=45}]
\node{简单机械}
child[concept color=green!25] { node {杠杆}
  child { node {五要素} }
  child { node {力臂} }
  child { node {最小力} }
  child { node {动态} }
}
child[concept color=orange!25] { node {滑轮}
  child { node {绳股数} }
  child { node {$F,s,v$关系} }
  child { node {水平拉物} }
  child { node {摩擦参与} }
}
child[concept color=purple!25] { node {机械效率}
  child { node {竖直提升} }
  child { node {水平牵引} }
  child { node {实验测量} }
  child { node {功率} }
}
child[concept color=red!25] { node {综合}
  child { node {浮力综合} }
  child { node {图像题} }
  child { node {极值题} }
  child { node {多过程} }
};
\end{tikzpicture}
\end{center}

\clearpage
\section{第一讲\texorpdfstring{\quad}{ }杠杆模型}

\subsection{知识框架}
\begin{itemize}[leftmargin=2.2em]
\item 杠杆五要素：支点、动力、阻力、动力臂、阻力臂。
\item 力臂本质：\textbf{支点到力的作用线的距离}，不是支点到作用点的距离。
\item 平衡条件：$F_1 l_1 = F_2 l_2$。
\item 最小力问题：本质是\textbf{在阻力矩一定时，使动力臂最大}。
\item 动态问题：常化为“一个量固定，另一个量变”，常用正弦、余弦投影。
\item 浮力综合：把浸在液体中的重物看成“\textbf{有效重力}” $G_{\text{效}}=G-F_b$。
\end{itemize}

\begin{center}
\begin{tikzpicture}[scale=0.95,>=Stealth]
\draw[thick] (0,0)--(6,0);
\fill (2,0) circle(2.2pt) node[below] {$O$};
\draw[fill=gray!35] (1.72,-0.18)--(2.28,-0.18)--(2,0.45)--cycle;
\draw[->,thick] (0.8,0)--(0.8,-1.5) node[below] {$F_1$};
\draw[->,thick] (5.1,0)--(5.1,1.6) node[above] {$F_2$};
\draw[dashed] (0.8,0.55)--(0.8,-1.55);
\draw[dashed] (5.1,-0.55)--(5.1,1.65);
\draw[decorate,decoration={brace,mirror,amplitude=5pt}] (0.8,-0.45)--(2,-0.45) node[midway,below=7pt] {$l_1$};
\draw[decorate,decoration={brace,amplitude=5pt}] (2,0.45)--(5.1,0.45) node[midway,above=7pt] {$l_2$};
\node at (3,-2.5) {关键：先找\textbf{作用线}，再作支点到作用线的距离};
\end{tikzpicture}
\end{center}

\subsection{核心模型}
\begin{modelbox}{本章核心模型}
\begin{enumerate}[label=\arabic*.]
\item \textbf{五要素模型：} 支点定住，作用线画出，力臂自然明朗。
\item \textbf{平衡模型：} 同向矩与反向矩相等，不急于代数，先判方向。
\item \textbf{最小力模型：} 阻力矩一定 $\Rightarrow$ 动力臂越大，动力越小。
\item \textbf{动态模型：} 把变化的力臂写成几何投影，如 $l=r\sin\theta$ 或 $r\cos\theta$。
\item \textbf{浮力替代模型：} 与液体接触的重物，先改写成 $G-F_b$ 再进杠杆。
\end{enumerate}
\end{modelbox}

\subsection{常见误区}
\begin{pitfallbox}{高失分点}
\begin{enumerate}[label=\arabic*.]
\item 把力臂误认为“支点到作用点的距离”。
\item 只会背“省力杠杆”，不会用力矩平衡定量求解。
\item 见到“最小力”就默认“沿杆方向拉”，这是错误的；最小力对应最大力臂，常常应\textbf{垂直于连线}。
\item 动态杠杆不先看哪个量定、哪个量变，直接代公式。
\item 浮力综合中忘了先求有效重力。
\end{enumerate}
\end{pitfallbox}

\subsection{典型例题}
\paragraph{例1}
轻质杠杆 AB 可绕 O 点转动，$OA=0.20\,\mathrm{m}$，$OB=0.60\,\mathrm{m}$。A 端挂重物
$G=45\,\mathrm{N}$，在 B 端竖直向下施加力 $F$，使杠杆水平平衡，求 $F$。

\Ans{$15\,\mathrm{N}$。}
\Sol{由杠杆平衡条件：
\[
G\cdot OA = F\cdot OB
\Rightarrow 45\times 0.20 = F\times 0.60
\]
解得 $F=15\,\mathrm{N}$。}
\Tip{杠杆题不要先问省力还是费力，先列力矩。}

\subsection{一题多解}
\paragraph{例2}
轻质杠杆 OA 长 $1.2\,\mathrm{m}$，O 为支点，A 端挂重物 $G=48\,\mathrm{N}$。在点 B 处施加力
$F$ 使杠杆平衡，且 $OB=0.8\,\mathrm{m}$。求：$F$ 的最小值及其方向。

\Ans{$F_{\min}=72\,\mathrm{N}$；方向应与 $OB$ 垂直。}

\Sol{\textbf{方法一：力矩法。}阻力矩一定为
\[
M=G\cdot OA =48\times 1.2=57.6\,\mathrm{N\cdot m}
\]
要使 $F$ 最小，应使动力臂最大。B 点到支点的最长距离就是 $OB=0.8\,\mathrm{m}$，当
$F\perp OB$ 时动力臂最大。于是
\[
F_{\min}=\frac{57.6}{0.8}=72\,\mathrm{N}.
\]}

\Sol{\textbf{方法二：几何法。}对固定点 B 来说，力的大小一定时，力矩最大当且仅当力的方向与
$OB$ 垂直。这是“定点施力最小力问题”的统一结论。本题所需阻力矩固定，因此最小力也在
$F\perp OB$ 时取得，大小同上为 $72\,\mathrm{N}$。}

\subsection{变式训练}

\subsubsection*{基础巩固}
\paragraph{基础1}
轻质杠杆 OA、OB 两侧长度分别为 $10\,\mathrm{cm}$ 和 $50\,\mathrm{cm}$，A 端挂 $20\,\mathrm{N}$ 重物，
在 B 端竖直向下施加力 $F$ 恰好平衡。求 $F$。

\paragraph{基础2}
一根长 $40\,\mathrm{cm}$ 的轻杆，以 O 点为支点，杆与水平方向成 $30^\circ$ 角。A 端受到竖直向下
$10\,\mathrm{N}$ 的力。求该力对支点 O 的力臂。

\paragraph{基础3}
下列工具分别属于省力杠杆还是费力杠杆：瓶起子、镊子、羊角锤拔钉。

\subsubsection*{中考提升}
\paragraph{提升1}
门把手等效为长 $0.60\,\mathrm{m}$ 的杠杆。若需要产生 $12\,\mathrm{N\cdot m}$ 的力矩才能转动门锁，
求最小作用力。

\paragraph{提升2}
轻质杠杆 OA 长 $0.80\,\mathrm{m}$，A 端挂重物 $G=20\,\mathrm{N}$。在点 B 施加水平拉力 $F$，
其中 $OB=0.40\,\mathrm{m}$。若杠杆与水平成 $\theta$ 角时静止，求 $F$ 关于 $\theta$ 的表达式，
并判断 $\theta$ 增大时 $F$ 怎样变化。

\paragraph{提升3}
水平轻质杠杆左端 A 处悬挂一重 $30\,\mathrm{N}$ 的物体，物体全部浸没在水中时所受浮力为
$10\,\mathrm{N}$。已知 $OA=20\,\mathrm{cm}$，$OB=60\,\mathrm{cm}$，在 B 端竖直向上施加力 $F$ 使杠杆水平平衡。
求 $F$。

\paragraph{提升4}
均匀杠杆 AB 长 $1.0\,\mathrm{m}$，重 $10\,\mathrm{N}$，支点 O 距 A 端 $0.20\,\mathrm{m}$。A 端挂
$20\,\mathrm{N}$ 重物。为使杠杆水平平衡，B 端应施加多大的力？方向向上还是向下？

\subsubsection*{竞赛拔高}
\paragraph{拔高1}
一根轻杆 AB 长 $1.2\,\mathrm{m}$，A 端挂 $10\,\mathrm{N}$，B 端挂 $30\,\mathrm{N}$，忽略杆重。
若要使杠杆水平平衡，支点 O 应放在距 A 端多远处？

\paragraph{拔高2}
轻质杠杆 OA 长 $1.0\,\mathrm{m}$，A 端挂重物 $G=30\,\mathrm{N}$。在点 B 施加力 $F$，其中
$OB=0.50\,\mathrm{m}$，且 $F$ 与杠杆的夹角为 $\theta$。若 $F$ 的大小不能超过 $80\,\mathrm{N}$，
求杠杆能平衡时 $\theta$ 的取值范围。

\paragraph{拔高3}
均匀杠杆 AB 长 $1.0\,\mathrm{m}$，重 $14\,\mathrm{N}$，支点 O 距 A 端 $0.25\,\mathrm{m}$。A 端悬挂一物体，
物重 $18\,\mathrm{N}$，浸没在水中后浮力为 $8\,\mathrm{N}$。求为了让杠杆水平平衡，B 端应施加多大的力？
方向向上还是向下？

\subsection{详细答案与解析}

\paragraph{基础1}
\Ans{$4\,\mathrm{N}$。}
\Sol{由平衡条件：
\[
20\times 10 = F\times 50
\Rightarrow F=4\,\mathrm{N}.
\]}

\paragraph{基础2}
\Ans{$20\sqrt{3}\,\mathrm{cm}$。}
\Sol{力竖直向下，所以力臂等于支点 O 到竖直作用线的水平距离。A 点相对 O 的水平距离为
\[
40\cos 30^\circ = 20\sqrt{3}\,\mathrm{cm}.
\]
故力臂为 $20\sqrt{3}\,\mathrm{cm}$。}

\paragraph{基础3}
\Ans{瓶起子：省力杠杆；镊子：费力杠杆；羊角锤拔钉：省力杠杆。}
\Sol{判断标准不是“看起来费不费劲”，而是看动力臂和阻力臂的大小关系。}

\paragraph{提升1}
\Ans{$20\,\mathrm{N}$。}
\Sol{最小力对应最大力臂，力臂最大取 $0.60\,\mathrm{m}$。于是
\[
F_{\min}=\frac{12}{0.60}=20\,\mathrm{N}.
\]}

\paragraph{提升2}
\Ans{$F=40\cot\theta\ \mathrm{N}$；$\theta$ 增大时 $F$ 减小。}
\Sol{A 端重力对支点的力臂为 $OA\cos\theta$；B 点水平力对支点的力臂为 $OB\sin\theta$。
由力矩平衡：
\[
F\cdot 0.40\sin\theta = 20\cdot 0.80\cos\theta
\]
故
\[
F=40\cot\theta.
\]
当 $\theta$ 增大时，$\cot\theta$ 减小，所以 $F$ 减小。}

\paragraph{提升3}
\Ans{$\dfrac{20}{3}\,\mathrm{N}$。}
\Sol{浸没后有效重力
\[
G_{\text{效}}=30-10=20\,\mathrm{N}.
\]
由平衡条件：
\[
20\times 20 = F\times 60
\Rightarrow F=\frac{20}{3}\,\mathrm{N}.
\]}

\paragraph{提升4}
\Ans{$1.25\,\mathrm{N}$，方向向下。}
\Sol{以 O 为支点。A 端重物在左侧，产生逆时针力矩：
\[
20\times 0.20 = 4.
\]
杆重作用点在中点，距 O 的距离为 $0.50-0.20=0.30\,\mathrm{m}$，产生顺时针力矩：
\[
10\times 0.30=3.
\]
还差 $1\,\mathrm{N\cdot m}$ 的顺时针力矩，所以 B 端应向下施力：
\[
F\times 0.80=1 \Rightarrow F=1.25\,\mathrm{N}.
\]}

\paragraph{拔高1}
\Ans{$0.90\,\mathrm{m}$。}
\Sol{设支点 O 距 A 端 $x$，则距 B 端为 $1.2-x$。由平衡条件：
\[
10x = 30(1.2-x)
\]
解得
\[
40x=36 \Rightarrow x=0.90\,\mathrm{m}.
\]
说明支点更靠近较重的一端。}

\paragraph{拔高2}
\Ans{$48.6^\circ \le \theta \le 131.4^\circ$。}
\Sol{力矩平衡：
\[
F\cdot 0.50\sin\theta = 30\times 1.0.
\]
所以
\[
F=\frac{60}{\sin\theta}.
\]
要使 $F\le 80$，必须有
\[
\frac{60}{\sin\theta}\le 80 \Rightarrow \sin\theta\ge 0.75.
\]
故
\[
\theta\in [\arcsin 0.75,\ 180^\circ-\arcsin 0.75]
\]
即
\[
48.6^\circ \le \theta \le 131.4^\circ.
\]}

\paragraph{拔高3}
\Ans{$\dfrac{4}{3}\,\mathrm{N}$，方向向上。}
\Sol{A 端物体浸没后有效重力
\[
18-8=10\,\mathrm{N}.
\]
其对 O 的力矩为
\[
10\times 0.25=2.5.
\]
杆重作用于杆中点，距 O 的距离为 $0.50-0.25=0.25\,\mathrm{m}$，力矩为
\[
14\times 0.25=3.5.
\]
右侧顺时针偏大，多出 $1.0\,\mathrm{N\cdot m}$，故 B 端应向上施力提供逆时针力矩：
\[
F\times 0.75=1.0 \Rightarrow F=\frac{4}{3}\,\mathrm{N}.
\]}

\clearpage
\section{第二讲\texorpdfstring{\quad}{ }滑轮与滑轮组}

\subsection{知识框架}
\begin{itemize}[leftmargin=2.2em]
\item 定滑轮：实质相当于等臂杠杆，主要改变力的方向。
\item 动滑轮：实质相当于动力臂为阻力臂 2 倍的杠杆。
\item 滑轮组第一步不是算，而是\textbf{数承担动滑轮的绳股数} $n$。
\item 理想情形：竖直提升时 $nF=G+G_{\text{动}}$；水平匀速拉物体时 $nF=f$。
\item 位移关系：$s=nh$；速度关系：$v_{\text{绳}}=nv_{\text{物}}$。
\item 同一滑轮组“越省力，越费距离”，但\textbf{不省功}。
\end{itemize}

\begin{center}
\begin{tikzpicture}[scale=0.95,>=Stealth]
\draw[thick] (0,4)--(6,4);
\draw[thick] (1.2,4)--(1.2,3.35);
\draw[thick] (4.8,4)--(4.8,3.35);

\draw[thick] (1.2,3.05) circle(0.28); % fixed pulley
\draw[thick] (4.0,1.65) circle(0.38); % movable pulley

\draw[thick] (0.4,4) -- (0.4,1.65);
\draw[thick] (0.4,1.65) -- (3.62,1.65);
\draw[thick] (4.38,1.65) -- (1.48,3.05);
\draw[thick] (0.92,3.05) -- (0.92,0.55);
\draw[->,thick] (0.92,0.55)--(0.92,-0.75) node[below] {$F$};

\draw[thick] (4.0,1.27)--(4.0,0.25);
\draw[fill=gray!25] (3.65,-0.55) rectangle (4.35,0.25);
\node at (4.0,-0.15) {$G$};

\node at (2.3,2.2) {$n=2$};
\node at (5.2,1.65) {动滑轮};
\node at (1.85,3.1) {定滑轮};
\end{tikzpicture}
\end{center}

\subsection{核心模型}
\begin{modelbox}{本章核心模型}
\begin{enumerate}[label=\arabic*.]
\item \textbf{数股模型：} 只数“直接承担动滑轮”的绳段，不数闲置段。
\item \textbf{受力模型：} 竖直提升写 $nF=\text{总重}$；水平匀速写 $nF=f$。
\item \textbf{位移速度模型：} 自由端走得远，物体走得近；自由端更快，物体更慢。
\item \textbf{功率模型：} $P=Fv_{\text{绳}}$，也可写成 $P_{\text{有用}}=f v_{\text{物}}$ 或 $Gv_{\text{物}}$。
\item \textbf{比较模型：} 同一物体同一速度，股数越多，拉力越小，自由端速度越大，功率理想情况下不变。
\end{enumerate}
\end{modelbox}

\begin{center}
\small
\begin{tabularx}{\textwidth}{c c c c X}
\toprule
绳股数 $n$ & 理想拉力 & 路程关系 & 速度关系 & 适用提醒\\
\midrule
1 & $F=G$ 或 $F=f$ & $s=h$ & $v_{\text{绳}}=v_{\text{物}}$ & 单个定滑轮，改方向不省力\\
2 & $F=\dfrac{G+G_{\text{动}}}{2}$ 或 $F=\dfrac{f}{2}$ & $s=2h$ & $v_{\text{绳}}=2v_{\text{物}}$ & 最常见\\
3 & $F=\dfrac{G+G_{\text{动}}}{3}$ 或 $F=\dfrac{f}{3}$ & $s=3h$ & $v_{\text{绳}}=3v_{\text{物}}$ & 压轴题常考\\
4 & $F=\dfrac{G+G_{\text{动}}}{4}$ 或 $F=\dfrac{f}{4}$ & $s=4h$ & $v_{\text{绳}}=4v_{\text{物}}$ & 易和功率混合考查\\
\bottomrule
\end{tabularx}
\end{center}

\subsection{常见误区}
\begin{pitfallbox}{高失分点}
\begin{enumerate}[label=\arabic*.]
\item 上来就套公式，没有先数 $n$。
\item 水平拉物体时还在写 $G=nF$，这是方向混乱。
\item 忘记自由端速度是物体速度的 $n$ 倍。
\item 以为动滑轮“省力也省功”。
\item 同一题中既求力又求功率时，速度对应错对象。
\end{enumerate}
\end{pitfallbox}

\subsection{典型例题}
\paragraph{例1}
某理想滑轮组承担动滑轮的绳股数为 4，提升重物 $G=140\,\mathrm{N}$，动滑轮总重
$20\,\mathrm{N}$。若物体上升 $0.30\,\mathrm{m}$，求拉力和自由端移动距离。

\Ans{$40\,\mathrm{N}$；$1.20\,\mathrm{m}$。}
\Sol{由竖直提升模型：
\[
4F=140+20=160 \Rightarrow F=40\,\mathrm{N}.
\]
又因为
\[
s=nh=4\times 0.30=1.20\,\mathrm{m}.
\]}

\subsection{一题多解}
\paragraph{例2}\mbox{}
\\
用三股绳滑轮组在水平地面上匀速拉动木箱，木箱所受摩擦力为 $90\,\mathrm{N}$，木箱速度为
$0.20\,\mathrm{m/s}$。求拉力 $F$ 与自由端功率 $P$。

\Ans{$F=30\,\mathrm{N}$；$P=18\,\mathrm{W}$。}

\Sol{\textbf{方法一：受力+速度。}
匀速运动时
\[
3F=f=90\,\mathrm{N}\Rightarrow F=30\,\mathrm{N}.
\]
自由端速度：
\[
v_{\text{绳}}=3v_{\text{物}}=3\times 0.20=0.60\,\mathrm{m/s}.
\]
故功率
\[
P=Fv_{\text{绳}}=30\times 0.60=18\,\mathrm{W}.
\]}

\Sol{\textbf{方法二：有用功率法。}
理想情况下输入功率等于有用功率：
\[
P=fv_{\text{物}}=90\times 0.20=18\,\mathrm{W}.
\]
再由 $P=Fv_{\text{绳}}$ 或 $3F=90$ 可得 $F=30\,\mathrm{N}$。}

\subsection{变式训练}

\subsubsection*{基础巩固}
\paragraph{基础1}
用理想定滑轮匀速提升 $60\,\mathrm{N}$ 的重物，求拉力大小，并说明是否省力。

\paragraph{基础2}
用理想动滑轮匀速提升重物 $100\,\mathrm{N}$，动滑轮重 $20\,\mathrm{N}$，求拉力。

\paragraph{基础3}
某滑轮组中承担动滑轮的绳股数为 3。当物体上升 $20\,\mathrm{cm}$ 时，自由端移动多远？
若物体速度为 $0.20\,\mathrm{m/s}$，自由端速度为多大？

\subsubsection*{中考提升}
\paragraph{提升1}
用三股绳滑轮组在粗糙水平面上匀速拉动物体，已知物体受到的摩擦力为
$36\,\mathrm{N}$，求拉力 $F$。

\paragraph{提升2}
在一套四股绳滑轮组中，自由端速度为 $0.80\,\mathrm{m/s}$。求物体速度；若持续
$10\,\mathrm{s}$，物体上升高度是多少？

\paragraph{提升3}
理想四股绳滑轮组提升重物 $140\,\mathrm{N}$，动滑轮重 $20\,\mathrm{N}$。求拉力；
若自由端向下拉 $1.20\,\mathrm{m}$，重物上升多少？

\paragraph{提升4}
用四股绳滑轮组水平匀速拉动木箱，木箱所受摩擦力为 $96\,\mathrm{N}$。
若自由端速度为 $1.0\,\mathrm{m/s}$，求拉力和拉力功率。

\subsubsection*{竞赛拔高}
\paragraph{拔高1}
某人站在随动滑轮一起上升的平台上，用手向下拉绳使自己和平台一起匀速上升。
已知“人+平台”总重为 $600\,\mathrm{N}$，忽略滑轮和绳重，求手的拉力。

\paragraph{拔高2}
用三股绳滑轮组在水平地面上拉木箱，木箱所受摩擦力为 $90\,\mathrm{N}$。
若自由端速度始终为 $0.60\,\mathrm{m/s}$，持续拉动 $10\,\mathrm{s}$。求：木箱位移、拉力做的总功、
有用功率。

\paragraph{拔高3}
用两股绳和三股绳两种理想滑轮组，分别匀速提升同一重物。已知重物重
$100\,\mathrm{N}$，动滑轮总重均为 $20\,\mathrm{N}$，重物上升速度均为 $0.10\,\mathrm{m/s}$。
分别求两种情况下的拉力和拉力功率，并说明“省力不省功”。

\subsection{详细答案与解析}

\paragraph{基础1}
\Ans{$60\,\mathrm{N}$；不省力，只改变方向。}
\Sol{定滑轮理想情况下动力臂和阻力臂相等，所以
\[
F=G=60\,\mathrm{N}.
\]}

\paragraph{基础2}
\Ans{$60\,\mathrm{N}$。}
\Sol{理想动滑轮有两股绳承担：
\[
2F=100+20=120 \Rightarrow F=60\,\mathrm{N}.
\]}

\paragraph{基础3}
\Ans{$60\,\mathrm{cm}$；$0.60\,\mathrm{m/s}$。}
\Sol{因为 $n=3$，
\[
s=nh=3\times 20\,\mathrm{cm}=60\,\mathrm{cm},
\]
\[
v_{\text{绳}}=nv_{\text{物}}=3\times 0.20=0.60\,\mathrm{m/s}.
\]}

\paragraph{提升1}
\Ans{$12\,\mathrm{N}$。}
\Sol{水平匀速：
\[
3F=f=36\,\mathrm{N}\Rightarrow F=12\,\mathrm{N}.
\]}

\paragraph{提升2}
\Ans{$0.20\,\mathrm{m/s}$；$2.0\,\mathrm{m}$。}
\Sol{四股绳关系：
\[
v_{\text{物}}=\frac{v_{\text{绳}}}{4}=\frac{0.80}{4}=0.20\,\mathrm{m/s}.
\]
故 10 s 内物体上升
\[
h=v_{\text{物}}t=0.20\times 10=2.0\,\mathrm{m}.
\]}

\paragraph{提升3}
\Ans{$40\,\mathrm{N}$；$0.30\,\mathrm{m}$。}
\Sol{理想情况下
\[
4F=140+20=160 \Rightarrow F=40\,\mathrm{N}.
\]
位移关系：
\[
h=\frac{s}{4}=\frac{1.20}{4}=0.30\,\mathrm{m}.
\]}

\paragraph{提升4}
\Ans{$24\,\mathrm{N}$；$24\,\mathrm{W}$。}
\Sol{匀速拉动时
\[
4F=96 \Rightarrow F=24\,\mathrm{N}.
\]
拉力功率：
\[
P=Fv_{\text{绳}}=24\times 1.0 =24\,\mathrm{W}.
\]}

\paragraph{拔高1}
\Ans{$200\,\mathrm{N}$。}
\Sol{把“人+平台+动滑轮”看作整体。外界向上的绳拉力有三段张力：两段作用在动滑轮上，
一段作用在人的手上，因此
\[
3F=600 \Rightarrow F=200\,\mathrm{N}.
\]
这是典型的“人拉自己”问题。}

\paragraph{拔高2}
\Ans{木箱位移 $2.0\,\mathrm{m}$；总功 $180\,\mathrm{J}$；有用功率 $18\,\mathrm{W}$。}
\Sol{由 $3F=90$ 得
\[
F=30\,\mathrm{N}.
\]
木箱速度
\[
v_{\text{物}}=\frac{0.60}{3}=0.20\,\mathrm{m/s}.
\]
10 s 内木箱位移
\[
x=v_{\text{物}}t=0.20\times 10=2.0\,\mathrm{m}.
\]
自由端位移
\[
s=0.60\times 10=6.0\,\mathrm{m}.
\]
总功
\[
W_{\text{总}}=Fs=30\times 6=180\,\mathrm{J}.
\]
有用功率
\[
P_{\text{有用}}=fv_{\text{物}}=90\times 0.20=18\,\mathrm{W}.
\]}

\paragraph{拔高3}
\Ans{两股绳：$F_2=60\,\mathrm{N}$，$P_2=12\,\mathrm{W}$；三股绳：$F_3=40\,\mathrm{N}$，
$P_3=12\,\mathrm{W}$。}
\Sol{总被提升重力均为
\[
100+20=120\,\mathrm{N}.
\]
两股绳：
\[
F_2=\frac{120}{2}=60\,\mathrm{N},\quad v_{\text{绳}}=2\times 0.10=0.20\,\mathrm{m/s}
\]
故
\[
P_2=F_2v_{\text{绳}}=60\times 0.20=12\,\mathrm{W}.
\]
三股绳：
\[
F_3=\frac{120}{3}=40\,\mathrm{N},\quad v_{\text{绳}}=3\times 0.10=0.30\,\mathrm{m/s}
\]
故
\[
P_3=40\times 0.30=12\,\mathrm{W}.
\]
可见绳股数越多，拉力越小，但自由端速度越大，因此理想情况下并不省功。}

\clearpage
\section{第三讲\texorpdfstring{\quad}{ }功、功率与机械效率}

\subsection{知识框架}
\begin{itemize}[leftmargin=2.2em]
\item 功：$W=Fs$，但要看\textbf{哪个力、哪个位移}。
\item 功率：$P=\dfrac{W}{t}=Fv$。
\item 机械效率：
\[
\eta=\frac{W_{\text{有用}}}{W_{\text{总}}}\times 100\%.
\]
\item 竖直提升模型：$W_{\text{有用}}=Gh$，$W_{\text{总}}=Fs$。
\item 水平拉物模型：$W_{\text{有用}}=fs_{\text{物}}$，$W_{\text{总}}=Fs_{\text{绳}}$。
\item 同一机械中，机械效率往往随负载增大而提高，但不是永远不变。
\end{itemize}

\begin{center}
\small
\begin{tabularx}{\textwidth}{>{\raggedright\arraybackslash}p{3.0cm} X}
\toprule
模型 & 核心表达式\\
\midrule
竖直提升型 &
\[
\eta=\frac{Gh}{Fs}=\frac{G}{nF}
\]
若只有动滑轮重导致额外功，则
\[
\eta=\frac{G}{G+G_{\text{动}}}
\]
\\[0.8em]
水平匀速牵引型 &
\[
\eta=\frac{fs_{\text{物}}}{Fs_{\text{绳}}}=\frac{f}{nF}
\]
\\[0.8em]
功率型 &
\[
P_{\text{总}}=Fv_{\text{绳}},\qquad
P_{\text{有用}}=Gv_{\text{物}} \text{ 或 } fv_{\text{物}}
\]
且
\[
\eta=\frac{P_{\text{有用}}}{P_{\text{总}}}
\]
\\
\bottomrule
\end{tabularx}
\end{center}

\subsection{核心模型}
\begin{modelbox}{本章核心模型}
\begin{enumerate}[label=\arabic*.]
\item \textbf{先分有用与总，再算效率。}
\item \textbf{先分“谁在动”，再写功率。}
\item \textbf{能用比值就别先求绝对量。} 很多效率题直接用 $\eta=\dfrac{G}{nF}$ 或 $\eta=\dfrac{f}{nF}$。
\item \textbf{实验题抓四量：} $G,h,F,s$，再固定变量做比较。
\item \textbf{图像/数据题抓常量：} 若 $2F-G$ 或 $nF-f$ 近似不变，说明机械的额外阻力近似稳定。
\end{enumerate}
\end{modelbox}

\subsection{常见误区}
\begin{pitfallbox}{高失分点}
\begin{enumerate}[label=\arabic*.]
\item 把机械效率写成“大于 1”的数。
\item 把动滑轮重算进有用功。
\item 位移写错对象：提升高度是物体走的，绳长是自由端走的。
\item 功率题中把 $v_{\text{绳}}$ 和 $v_{\text{物}}$ 混用。
\item 实验结论不写“在同一机械、其他条件相同时”。
\end{enumerate}
\end{pitfallbox}

\subsection{典型例题}
\paragraph{例1}
用两股绳滑轮组提升重物 $G=200\,\mathrm{N}$，物体上升高度 $h=2.0\,\mathrm{m}$，
拉力 $F=120\,\mathrm{N}$。求有用功、总功和机械效率。

\Ans{$W_{\text{有用}}=400\,\mathrm{J}$；$W_{\text{总}}=480\,\mathrm{J}$；$\eta=83.3\%$。}
\Sol{两股绳时自由端位移
\[
s=2h=4.0\,\mathrm{m}.
\]
故
\[
W_{\text{有用}}=Gh=200\times 2=400\,\mathrm{J},
\]
\[
W_{\text{总}}=Fs=120\times 4=480\,\mathrm{J}.
\]
机械效率
\[
\eta=\frac{400}{480}=83.3\%.
\]}

\subsection{一题多解}
\paragraph{例2}
某四股绳滑轮组匀速提升重物，已知拉力 $F=25\,\mathrm{N}$，自由端速度
$v_{\text{绳}}=0.80\,\mathrm{m/s}$，被提升重物重 $G=70\,\mathrm{N}$。求机械效率。

\Ans{$70\%$。}

\Sol{\textbf{方法一：功率法。}
由速度关系：
\[
v_{\text{物}}=\frac{0.80}{4}=0.20\,\mathrm{m/s}.
\]
总功率
\[
P_{\text{总}}=Fv_{\text{绳}}=25\times 0.80=20\,\mathrm{W}.
\]
有用功率
\[
P_{\text{有用}}=Gv_{\text{物}}=70\times 0.20=14\,\mathrm{W}.
\]
故
\[
\eta=\frac{14}{20}=70\%.
\]}

\Sol{\textbf{方法二：比值法。}
由
\[
\eta=\frac{G}{nF}
\]
直接得
\[
\eta=\frac{70}{4\times 25}=\frac{70}{100}=70\%.
\]
竞赛和压轴中，能直接用比值法通常更快。}

\subsection{变式训练}

\subsubsection*{基础巩固}
\paragraph{基础1}
用两股绳滑轮组提升 $200\,\mathrm{N}$ 的重物 $2.0\,\mathrm{m}$，拉力为
$120\,\mathrm{N}$。求有用功、总功和机械效率。

\paragraph{基础2}
若自由端拉力为 $50\,\mathrm{N}$，自由端速度为 $0.60\,\mathrm{m/s}$，
求拉力功率。

\paragraph{基础3}
某机械做总功 $500\,\mathrm{J}$，其中有用功为 $350\,\mathrm{J}$。
求额外功和机械效率。

\subsubsection*{中考提升}
\paragraph{提升1}
理想两股绳滑轮组提升重物 $180\,\mathrm{N}$，动滑轮重 $20\,\mathrm{N}$。
求机械效率。

\paragraph{提升2}
某三股绳滑轮组在水平面上匀速拉动物体，物体受到的摩擦力为
$120\,\mathrm{N}$，拉力为 $50\,\mathrm{N}$。求机械效率。

\paragraph{提升3}
测滑轮组机械效率实验中，某同学记录数据如下：\\
第一次：$G=2.0\,\mathrm{N}$，$h=0.10\,\mathrm{m}$，$F=1.4\,\mathrm{N}$，$s=0.20\,\mathrm{m}$；\\
第二次：$G=3.0\,\mathrm{N}$，$h=0.10\,\mathrm{m}$，$F=1.9\,\mathrm{N}$，$s=0.20\,\mathrm{m}$。\\
求第二次的机械效率，并比较两次实验，写出一个结论。

\paragraph{提升4}
某四股绳滑轮组提升重物时，自由端拉力 $F=25\,\mathrm{N}$，自由端速度
$v_{\text{绳}}=0.80\,\mathrm{m/s}$，重物重 $70\,\mathrm{N}$。求物体速度、有用功率、总功率和机械效率。

\subsubsection*{竞赛拔高}
\paragraph{拔高1}
某滑轮组机械效率为 $75\%$。若忽略绳重和摩擦，且机械效率只由动滑轮重引起，
动滑轮重为 $30\,\mathrm{N}$。求所提升物体重力。

\paragraph{拔高2}
某两股绳滑轮组在水平地面上匀速拉动物体，拉力为 $40\,\mathrm{N}$，机械效率为
$75\%$。若自由端在 $10\,\mathrm{s}$ 内移动 $4.0\,\mathrm{m}$，求物体所受摩擦力、有用功和总功。

\paragraph{拔高3}
某同一滑轮组（两股绳）在不同负载下测得数据：
\[
(G,F)=(1.0\,\mathrm{N},0.8\,\mathrm{N}),\ (2.0\,\mathrm{N},1.3\,\mathrm{N}),\ (3.0\,\mathrm{N},1.8\,\mathrm{N})
\]
且每次都有 $s=2h$。试判断该机械的“额外阻力（等效）”是否近似不变，
并求其大小；再求第三次实验的机械效率。

\subsection{详细答案与解析}

\paragraph{基础1}
\Ans{$400\,\mathrm{J}$；$480\,\mathrm{J}$；$83.3\%$。}
\Sol{同例1：
\[
W_{\text{有用}}=200\times 2=400\,\mathrm{J},
\]
\[
s=2h=4\,\mathrm{m},\quad W_{\text{总}}=120\times 4=480\,\mathrm{J},
\]
\[
\eta=\frac{400}{480}=83.3\%.
\]}

\paragraph{基础2}
\Ans{$30\,\mathrm{W}$。}
\Sol{由
\[
P=Fv=50\times 0.60=30\,\mathrm{W}.
\]}

\paragraph{基础3}
\Ans{$150\,\mathrm{J}$；$70\%$。}
\Sol{额外功
\[
W_{\text{额}}=500-350=150\,\mathrm{J}.
\]
机械效率
\[
\eta=\frac{350}{500}=70\%.
\]}

\paragraph{提升1}
\Ans{$90\%$。}
\Sol{理想且仅由动滑轮重导致额外功：
\[
\eta=\frac{G}{G+G_{\text{动}}}=\frac{180}{180+20}=90\%.
\]}

\paragraph{提升2}
\Ans{$80\%$。}
\Sol{水平匀速拉物体时
\[
\eta=\frac{f}{nF}=\frac{120}{3\times 50}=80\%.
\]}

\paragraph{提升3}
\Ans{第二次机械效率约为 $78.9\%$；结论：在同一滑轮组、其他条件相同时，
提升物重越大，机械效率通常越高。}
\Sol{第二次：
\[
W_{\text{有用}}=Gh=3.0\times 0.10=0.30\,\mathrm{J},
\]
\[
W_{\text{总}}=Fs=1.9\times 0.20=0.38\,\mathrm{J}.
\]
故
\[
\eta=\frac{0.30}{0.38}\approx 78.9\%.
\]
第一次效率为
\[
\frac{2.0\times 0.10}{1.4\times 0.20}=\frac{0.20}{0.28}\approx 71.4\%.
\]
第二次更高。}

\paragraph{提升4}
\Ans{$0.20\,\mathrm{m/s}$；$14\,\mathrm{W}$；$20\,\mathrm{W}$；$70\%$。}
\Sol{由四股绳关系：
\[
v_{\text{物}}=\frac{0.80}{4}=0.20\,\mathrm{m/s}.
\]
有用功率：
\[
P_{\text{有用}}=Gv_{\text{物}}=70\times 0.20=14\,\mathrm{W}.
\]
总功率：
\[
P_{\text{总}}=Fv_{\text{绳}}=25\times 0.80=20\,\mathrm{W}.
\]
故
\[
\eta=\frac{14}{20}=70\%.
\]}

\paragraph{拔高1}
\Ans{$90\,\mathrm{N}$。}
\Sol{由
\[
\eta=\frac{G}{G+30}=75\%=0.75
\]
得
\[
G=0.75G+22.5
\Rightarrow 0.25G=22.5
\Rightarrow G=90\,\mathrm{N}.
\]}

\paragraph{拔高2}
\Ans{$60\,\mathrm{N}$；$120\,\mathrm{J}$；$160\,\mathrm{J}$。}
\Sol{两股绳时
\[
\eta=\frac{f}{2F}=75\%
\Rightarrow \frac{f}{80}=0.75
\Rightarrow f=60\,\mathrm{N}.
\]
自由端 10 s 位移 $4.0\,\mathrm{m}$，则物体位移
\[
x=\frac{4.0}{2}=2.0\,\mathrm{m}.
\]
有用功
\[
W_{\text{有用}}=fx=60\times 2=120\,\mathrm{J}.
\]
总功
\[
W_{\text{总}}=Fs=40\times 4=160\,\mathrm{J}.
\]}

\paragraph{拔高3}
\Ans{额外阻力近似不变，大小为 $0.6\,\mathrm{N}$；第三次机械效率约为 $83.3\%$。}
\Sol{两股绳时若将额外阻力等效为 $R$，则
\[
2F=G+R.
\]
分别代入三组数据：
\[
R=2\times 0.8-1.0=0.6\,\mathrm{N},
\]
\[
R=2\times 1.3-2.0=0.6\,\mathrm{N},
\]
\[
R=2\times 1.8-3.0=0.6\,\mathrm{N}.
\]
可见额外阻力近似恒定，为 $0.6\,\mathrm{N}$。第三次效率：
\[
\eta=\frac{G}{2F}=\frac{3.0}{3.6}=83.3\%.
\]}

\clearpage
\section{第四讲\texorpdfstring{\quad}{ }中考压轴与竞赛综合}

\subsection{知识框架}
\begin{itemize}[leftmargin=2.2em]
\item 图像题：先认横纵坐标，再判断线性关系的斜率和截距代表什么。
\item 极值题：本质往往是“力矩一定求最小力”或“受力表达式中的某一项最大/最小”。
\item 多过程题：一个过程一张小草图，一段位移一套方程，不许整题只列一式。
\item 杠杆+浮力：有效重力变化 $\Rightarrow$ 所需力线性变化。
\item 滑轮+摩擦+浮力：通常写成
\[
nF=f=\mu N=\mu (G-F_b)
\]
从而得出 $F$ 与浮力（或深度）的线性关系。
\item 实验设计题：一定写“控制变量”“测量量”“结论句式”。
\end{itemize}

\begin{modelbox}{本章核心模型}
\begin{enumerate}[label=\arabic*.]
\item \textbf{线性图像模型：} 只要 $F_b$ 与深度成正比、摩擦与支持力成正比，则拉力对深度通常是一条直线。
\item \textbf{多过程模型：} 分段求解后再求总功、平均力、总效率。
\item \textbf{组合机械模型：} 先把一个简单机械“等效”为某点上的力，再接第二个简单机械。
\item \textbf{反演模型：} 由图像的截距和斜率反推 $\mu$、体积、底面积、绳股数等参数。
\end{enumerate}
\end{modelbox}

\subsection{常见误区}
\begin{pitfallbox}{高失分点}
\begin{enumerate}[label=\arabic*.]
\item 图像题只盯数值，不解释斜率和截距。
\item 多过程题中“每段的位移对象”混乱。
\item 浮力变了，摩擦却还按原来支持力算。
\item 组合机械里忘记先做“中间量传递”。
\item 开放实验题只写器材，不写控制量与结论句。
\end{enumerate}
\end{pitfallbox}

\subsection{典型例题}
\paragraph{例1}
重为 $40\,\mathrm{N}$ 的木块置于粗糙水平面上，用两股绳滑轮组拉动。若木块浸入水后所受浮力为
$10\,\mathrm{N}$，且与地面的动摩擦因数可按 $f=\mu N$ 计算，其中 $\mu=0.4$。求木块匀速运动时的拉力。

\Ans{$6\,\mathrm{N}$。}
\Sol{浸在水中但仍压在地面上时，支持力
\[
N=G-F_b=40-10=30\,\mathrm{N}.
\]
故摩擦力
\[
f=\mu N=0.4\times 30=12\,\mathrm{N}.
\]
两股绳匀速拉动时
\[
2F=f=12\,\mathrm{N}\Rightarrow F=6\,\mathrm{N}.
\]}

\subsection{一题多解}
\paragraph{例2}
杠杆右端 B 通过挂钩悬挂一个理想动滑轮，动滑轮由两股绳承担。左端 A 悬挂一浸没在水中的物体，
物重为 $60\,\mathrm{N}$，浮力为 $12\,\mathrm{N}$。已知 $OA=0.20\,\mathrm{m}$，$OB=0.60\,\mathrm{m}$，
杠杆水平平衡。求自由端拉力 $F$。

\Ans{$8\,\mathrm{N}$。}

\Sol{\textbf{方法一：先杠杆，后滑轮。}
A 端有效重力为
\[
60-12=48\,\mathrm{N}.
\]
设动滑轮挂钩对杠杆 B 点的向下拉力为 $T_B$，由杠杆平衡：
\[
48\times 0.20 = T_B\times 0.60
\Rightarrow T_B=16\,\mathrm{N}.
\]
理想动滑轮由两股绳承担，所以
\[
2F=T_B=16 \Rightarrow F=8\,\mathrm{N}.
\]}

\Sol{\textbf{方法二：机械增益连乘。}
先看杠杆：左端 $48\,\mathrm{N}$ 通过臂长比
\[
\frac{OA}{OB}=\frac{1}{3}
\]
传到 B 点，故 B 点等效受力为
\[
48\times \frac{1}{3}=16\,\mathrm{N}.
\]
再看动滑轮：自由端拉力是挂钩载荷的一半，所以
\[
F=\frac{16}{2}=8\,\mathrm{N}.
\]
这类题的关键不是硬算，而是“一级机械的输出，变成下一级机械的输入”。}

\subsection{变式训练}

\subsubsection*{基础巩固}
\paragraph{基础1}
某杠杆所受阻力矩恒为 $9\,\mathrm{N\cdot m}$。当动力臂取 $0.30\,\mathrm{m}$ 时，求动力大小；
若拉力最大不能超过 $45\,\mathrm{N}$，则动力臂至少应为多大？

\paragraph{基础2}
重为 $40\,\mathrm{N}$ 的木块放在粗糙水平面上，用两股绳滑轮组匀速拉动。
木块浸入水后所受浮力为 $10\,\mathrm{N}$，摩擦满足 $f=0.4N$。求拉力。

\paragraph{基础3}
某三股绳滑轮组在水平面上匀速拉动物体。前 $2.0\,\mathrm{m}$ 运动中摩擦力为
$60\,\mathrm{N}$，后 $3.0\,\mathrm{m}$ 运动中摩擦力为 $90\,\mathrm{N}$。求两段中的拉力和整个过程拉力做的总功。

\subsubsection*{中考提升}
\paragraph{提升1}
某两股绳滑轮组拉动粗糙水平面上的木块，木块重 $40\,\mathrm{N}$，
摩擦满足 $f=0.4N$。若木块浸入液体后浮力随深度 $h$（单位：m）满足
$F_b=20h\,\mathrm{N}$，且在研究范围内始终未离开地面。求拉力 $F$ 关于深度 $h$ 的表达式，
并求当 $F=5\,\mathrm{N}$ 时的深度。

\paragraph{提升2}
一水平杠杆左端悬挂重物 $20\,\mathrm{N}$，浸入液体时浮力随浸入深度 $h$（单位：cm）
线性变化：$F_b=0.8h\,\mathrm{N}$，且 $0\le h\le 10$。已知左臂长 $20\,\mathrm{cm}$，
右臂长 $60\,\mathrm{cm}$。若右端竖直向上施力 $F$ 使杠杆平衡，求 $F$ 关于 $h$ 的表达式，
并求 $h=10\,\mathrm{cm}$ 时的 $F$。

\paragraph{提升3}
某两股绳滑轮组拉动木块，自由端速度恒为 $0.60\,\mathrm{m/s}$。若木块所受摩擦力
由 $30\,\mathrm{N}$ 降为 $18\,\mathrm{N}$，求前后两种状态下的拉力和拉力功率。

\paragraph{提升4}
现有弹簧测力计、刻度尺、钩码、两个相同滑轮、细绳和铁架台。请设计一个实验：
比较“绳股数不同”对滑轮组机械效率的影响。要求写出：控制变量、应测物理量、主要步骤、
判断结论的表达式。

\subsubsection*{竞赛拔高}
\paragraph{拔高1}
某理想三股绳滑轮组中，自由端拉力 $F$ 随物体位移 $x$ 的关系为一直线：
当 $x=0$ 时，$F=10\,\mathrm{N}$；当 $x=2.0\,\mathrm{m}$ 时，$F=6\,\mathrm{N}$。
求物体在这 $2.0\,\mathrm{m}$ 位移过程中，自由端拉力做的总功。

\paragraph{拔高2}
某两股绳滑轮组拉动木块在粗糙水平面上匀速运动。实验测得自由端拉力
$F$ 与浸入深度 $h$ 的图像是一条直线：当 $h=0$ 时，$F=12\,\mathrm{N}$；
当 $h=0.20\,\mathrm{m}$ 时，$F=4\,\mathrm{N}$。已知木块重 $60\,\mathrm{N}$，
且始终未离地，摩擦满足 $f=\mu N$，浮力满足 $F_b=\rho gSh$。
求摩擦因数 $\mu$ 和木块底面积 $S$。

\paragraph{拔高3}
重为 $80\,\mathrm{N}$ 的木块放在粗糙水平面上，用两股绳滑轮组匀速拉动。
若木块在连续三个阶段中所受浮力分别为 $0\,\mathrm{N}$、$20\,\mathrm{N}$、$40\,\mathrm{N}$，
且每个阶段木块都前进 $1.0\,\mathrm{m}$。已知摩擦满足 $f=0.3N$。
求三阶段的拉力、整个过程拉力做的总功和平均拉力。

\subsection{详细答案与解析}

\paragraph{基础1}
\Ans{$30\,\mathrm{N}$；$0.20\,\mathrm{m}$。}
\Sol{由
\[
F=\frac{M}{l}
\]
得
\[
F=\frac{9}{0.30}=30\,\mathrm{N}.
\]
若 $F\le 45\,\mathrm{N}$，则
\[
l\ge \frac{9}{45}=0.20\,\mathrm{m}.
\]}

\paragraph{基础2}
\Ans{$6\,\mathrm{N}$。}
\Sol{支持力
\[
N=40-10=30\,\mathrm{N},
\]
摩擦
\[
f=0.4\times 30=12\,\mathrm{N}.
\]
两股绳匀速拉动：
\[
2F=12 \Rightarrow F=6\,\mathrm{N}.
\]}

\paragraph{基础3}
\Ans{前段 $20\,\mathrm{N}$，后段 $30\,\mathrm{N}$，总功 $390\,\mathrm{J}$。}
\Sol{前段：
\[
3F_1=60 \Rightarrow F_1=20\,\mathrm{N}.
\]
后段：
\[
3F_2=90 \Rightarrow F_2=30\,\mathrm{N}.
\]
自由端位移分别为 $6\,\mathrm{m}$ 和 $9\,\mathrm{m}$，总功
\[
W=20\times 6 + 30\times 9=390\,\mathrm{J}.
\]}

\paragraph{提升1}
\Ans{$F=8-4h\ (\mathrm{N})$；当 $F=5\,\mathrm{N}$ 时，$h=0.75\,\mathrm{m}$。}
\Sol{木块仍压在地面上时
\[
N=40-F_b=40-20h.
\]
摩擦
\[
f=0.4N=0.4(40-20h)=16-8h.
\]
两股绳匀速拉动：
\[
2F=f \Rightarrow F=8-4h.
\]
令 $F=5$，
\[
8-4h=5 \Rightarrow h=0.75\,\mathrm{m}.
\]}

\paragraph{提升2}
\Ans{$F=\dfrac{20-0.8h}{3}\ (\mathrm{N})$；当 $h=10\,\mathrm{cm}$ 时，$F=4\,\mathrm{N}$。}
\Sol{左端有效重力
\[
G_{\text{效}}=20-0.8h.
\]
杠杆平衡：
\[
F\times 60 = (20-0.8h)\times 20.
\]
故
\[
F=\frac{20-0.8h}{3}.
\]
当 $h=10$ 时，
\[
F=\frac{20-8}{3}=4\,\mathrm{N}.
\]}

\paragraph{提升3}
\Ans{前：$F=15\,\mathrm{N}$，$P=9\,\mathrm{W}$；后：$F=9\,\mathrm{N}$，$P=5.4\,\mathrm{W}$。}
\Sol{两股绳匀速拉动：
\[
2F_1=30 \Rightarrow F_1=15\,\mathrm{N},
\qquad
2F_2=18 \Rightarrow F_2=9\,\mathrm{N}.
\]
拉力功率：
\[
P_1=F_1v_{\text{绳}}=15\times 0.60=9\,\mathrm{W},
\]
\[
P_2=9\times 0.60=5.4\,\mathrm{W}.
\]}

\paragraph{提升4}
\Ans{实验设计略，参考要点如下。}
\Sol{控制变量：使用同一组滑轮、相同提升高度、尽量保持匀速拉动，比较不同绳股数。
应测量：钩码重力 $G$、提升高度 $h$、拉力 $F$、自由端位移 $s$。
主要步骤：组装两种不同绳股数的滑轮组；分别匀速提升同一组钩码；记录 $G,h,F,s$；
计算
\[
\eta=\frac{Gh}{Fs}.
\]
结论句式：在同一机械、重物和提升高度相同时，比较不同绳股数对应的机械效率大小，
判断其影响。}

\paragraph{拔高1}
\Ans{$48\,\mathrm{J}$。}
\Sol{因为 $F$ 随 $x$ 线性变化，所以平均拉力
\[
\bar F=\frac{10+6}{2}=8\,\mathrm{N}.
\]
物体位移 $2.0\,\mathrm{m}$，三股绳对应自由端位移
\[
s=3\times 2.0=6.0\,\mathrm{m}.
\]
故总功
\[
W=\bar F\, s = 8\times 6=48\,\mathrm{J}.
\]
图像题常用“平均值乘总路程”的想法。}

\paragraph{拔高2}
\Ans{$\mu=0.4$；$S=0.020\,\mathrm{m^2}$。}
\Sol{两股绳匀速拉动：
\[
2F=f=\mu N=\mu (G-F_b).
\]
当 $h=0$ 时，$F=12\,\mathrm{N}$，此时 $F_b=0$，所以
\[
2\times 12=\mu \times 60 \Rightarrow \mu=0.4.
\]
又
\[
F=\frac{\mu(60-\rho gSh)}{2}.
\]
由图像斜率：
\[
\left|\frac{\Delta F}{\Delta h}\right|
=\frac{12-4}{0.20}=40\,\mathrm{N/m}.
\]
而理论斜率满足
\[
\frac{\mu \rho g S}{2}=40.
\]
代入 $\mu=0.4$、$\rho g=10^4\,\mathrm{N/m^3}$：
\[
\frac{0.4\times 10^4 \times S}{2}=40
\Rightarrow 2000S=40
\Rightarrow S=0.020\,\mathrm{m^2}.
\]}

\paragraph{拔高3}
\Ans{三阶段拉力分别为 $12\,\mathrm{N}$、$9\,\mathrm{N}$、$6\,\mathrm{N}$；
总功 $54\,\mathrm{J}$；平均拉力 $9\,\mathrm{N}$。}
\Sol{三阶段支持力分别为
\[
N_1=80,\quad N_2=80-20=60,\quad N_3=80-40=40.
\]
摩擦力分别为
\[
f_1=0.3\times 80=24\,\mathrm{N},
\]
\[
f_2=0.3\times 60=18\,\mathrm{N},
\]
\[
f_3=0.3\times 40=12\,\mathrm{N}.
\]
两股绳匀速拉动得
\[
F_1=12\,\mathrm{N},\quad F_2=9\,\mathrm{N},\quad F_3=6\,\mathrm{N}.
\]
每阶段木块前进 $1.0\,\mathrm{m}$，自由端各移动 $2.0\,\mathrm{m}$，总功
\[
W=12\times 2+9\times 2+6\times 2=54\,\mathrm{J}.
\]
总自由端位移
\[
s_{\text{总}}=2+2+2=6\,\mathrm{m},
\]
故平均拉力
\[
\bar F=\frac{W}{s_{\text{总}}}=\frac{54}{6}=9\,\mathrm{N}.
\]}

\section*{课堂使用建议}
\addcontentsline{toc}{section}{课堂使用建议}
\begin{itemize}[leftmargin=2.2em]
\item 第一轮：讲模型，不追求题量，重点是“会说出为什么这样列式”。
\item 第二轮：把每章提升题单独限时，训练中考压轴节奏。
\item 第三轮：拔高题只要求“会拆模型、敢下第一步”，不要一开始就追求满分。
\item 若学生基础较弱：先删去第四讲拔高2、拔高3，再做第一讲和第三讲的提升题。
\item 若学生冲竞赛：每道拔高题要求再写出一种替代思路。
\end{itemize}

\end{document}


